\chapter{Build: Tick-to-Trade, Measured}\label{ch:ll:build-tick-to-trade-measured}

A packet leaves the simulated exchange on two lines at once. A few microseconds later the \omterm{def:ll:the-feed-handler:handler}{feed handler} has it, a microsecond
after that the strategy has decided, and an order leaves for the venue through the \omterm{def:ll:pre-trade-risk-and-kill-switches:gate}{risk gate} and the \omterm{def:ll:order-gateway-and-order-management:gateway}{order gateway}. Every one
of those steps has been built and measured in isolation in this book; this chapter puts them on one path, on four threads of
one machine, and times the whole thing with a stamp at every boundary. The result is the \omterm{def:ll:where-latency-comes-from:budget}{latency budget} of chapter~1 with the
guesses replaced by measurements: which stage costs what, at the median and in the tail, and where a real server would have to
be different. The answer on this laptop is sobering in one place and reassuring in another. The code the book has written
takes a few hundred nanoseconds; the kernel, the network stack and a shared machine take the rest, and the rest is most of it.

\section{Assembling the path}

The path is the book's components in order (\cref{fig:ll:build-tick-to-trade-measured:path}). An \emph{exchange} thread replays
the recorded line: the small-tick day of chapter~19, 6\,000 events in 1\,200 MoldUDP64 packets, sent over UDP on the loopback
interface on line A and, five microseconds later, on line B, at twenty times its recorded spacing (a packet every
\qty{100}{\micro\second}). A \emph{feed} thread receives both lines and runs the \omterm{def:ll:the-feed-handler:handler}{feed handler} of chapter~18, which arbitrates
them and publishes each event on a ring (chapter~12). An \emph{engine} thread takes each event off the ring and runs
\texttt{firm.ticktotrade}'s \texttt{Path}: the strategy engine of chapter~20, with its \omterm{def:ll:the-order-book-builder:builder}{book builder} (chapter~19), then, for each
order it decides, the \omterm{def:ll:pre-trade-risk-and-kill-switches:gate}{risk gate} of chapter~22 and the \omterm{def:ll:order-gateway-and-order-management:gateway}{order gateway} of chapter~21, which writes the venue's binary message
(chapter~16) and sends it on a TCP session. A \emph{venue} thread receives the session's frames. Each thread is pinned to its
own CPU (chapter~4), and nothing on the path allocates after warm-up (chapter~6): the harness counts, and finds zero in every
run.

\begin{omfigure}
\begin{tikzpicture}[font=\scriptsize,>=Stealth,box/.style={draw,rounded corners=2pt,minimum height=0.8cm,align=center,
  inner sep=3pt,fill=omDef!12},pt/.style={circle,fill=omThm,inner sep=1.4pt}]
\node[box,fill=gray!12] (ex) at (0,0) {exchange\\ thread};
\node[box] (fd) at (3.3,0) {feed handler\\ (feed thread)};
\node[box,fill=omMeth!15] (rg) at (5.6,0) {ring};
\node[box,minimum width=4.2cm] (en) at (8.8,0) {engine thread: book, strategy,\\ risk gate, gateway};
\node[box,fill=gray!12] (vn) at (12.4,0) {venue\\ thread};
\draw[->] (ex) -- node[above=1pt,font=\tiny]{UDP A, B} (fd);
\draw[->] (fd) -- (rg);
\draw[->] (rg) -- (en);
\draw[->] (en) -- node[above]{TCP} (vn);
\foreach \x/\l in {0.75/exch,1.95/recv,4.2/pub,6.15/pop,11.35/sent,11.6/venue} {\node[pt] at (\x,-0.62) {};}
\node[below,align=center] at (0.75,-0.7) {exch};
\node[below,align=center] at (1.95,-0.7) {recv};
\node[below,align=center] at (4.2,-0.7) {pub};
\node[below,align=center] at (6.15,-0.7) {pop};
\node[below,align=center] at (8.8,-0.7) {book, decide, risk, encode};
\node[below,align=center] at (11.9,-0.7) {sent, venue};
\foreach \x in {7.3,8.2,9.2,10.2} {\node[pt] at (\x,-0.62) {};}
\draw[gray] (0.3,-0.62) -- (12.8,-0.62);
\end{tikzpicture}

\omcaption{The assembled path and its \omterm{def:ll:build-tick-to-trade-measured:points}{instrumentation points}. Four threads on four CPUs of one laptop; every order carries ten
\omterm{def:ll:cpu-microarchitecture:tsc}{time-stamp counter} readings, from the exchange's send of the packet that caused it to the venue's receipt of the
order.}\label{fig:ll:build-tick-to-trade-measured:path}
\end{omfigure}

The path is checked before it is timed. The C++ test replays the same events offline through the same \texttt{Path} and
obtains 1\,050 orders (new orders, replaces and cancels) whose hash is committed; the Rust twin, assembled from the Rust \omterm{def:ll:the-feed-handler:handler}{feed
handler}, engine, \omterm{def:ll:pre-trade-risk-and-kill-switches:gate}{risk gate} and gateway, produces the same hash; and the live harness, with its threads, sockets and
arbitration between two lines, produces it again, in every run. What is timed is therefore the path that computes the right
orders.

\omcode{code/firm/ticktotrade/cpp/firm_ticktotrade.hpp}{82}{91}{The engine thread's work for one market event: the engine
decides, then each order passes the gate and the gateway.}\label{lst:ll:build-tick-to-trade-measured:path}

\section{Instrumentation points and stage timestamps}

\begin{definition}[Instrumentation point, latency attribution]\label{def:ll:build-tick-to-trade-measured:points}
An \emph{instrumentation point}\index{instrumentation point} is a place in a path where the time is read and recorded with the
identity of the message passing it. \emph{Latency attribution}\index{latency attribution} divides a message's end-to-end latency
into the intervals between consecutive instrumentation points, so that each part is charged to the stage that spent it.
\end{definition}

Every point reads the \omterm{def:ll:cpu-microarchitecture:tsc}{time-stamp counter} (chapter~5), which all cores of this machine share, so that a reading on the feed
thread and one on the engine thread are comparable. The ten points of \cref{fig:ll:build-tick-to-trade-measured:path} give eight
stages: \emph{receive}, from the exchange's \texttt{sendto} to the feed thread's \texttt{recv}, the kernel's UDP path on the
loopback interface; \emph{decode}, the handler's arbitration and normalisation until the event is on the ring; \emph{ring},
the wait on the ring; \emph{book}, \emph{strategy}, \emph{risk} and \emph{gateway} on the engine thread; and \emph{transmit},
from the written message to the venue thread's receipt, the kernel's TCP path. The stamps travel with the event (in the ring's
message) and with the order (in a preallocated record), and the exchange's and venue's stamps are joined afterwards by packet
sequence number and client order identifier (\cref{lst:ll:build-tick-to-trade-measured:loop}). Two traps shaped the design.
Reading the clock is not free: each point costs some \qty{11}{\nano\second} with \texttt{rdtscp} (chapter~5), so the points
are few and none is inside a loop. And
when an event causes two orders, the second waits for the first to be sent, a few microseconds that the naive attribution
charges to its risk check; the stages are therefore attributed on the first order of each event, and the end-to-end figures on
every order.

\omcode{code/firm/ticktotrade/cpp/firm_ticktotrade_live.hpp}{162}{187}{The engine thread: take an event off the ring, stamp it,
run the path, send each order and stamp it.}\label{lst:ll:build-tick-to-trade-measured:loop}

\section{The budget against the measurement}

\Cref{fig:ll:build-tick-to-trade-measured:stages} gives the stages over ten runs, 5\,250 first orders (10\,500 in all), on this
laptop with no isolated cores, the machine otherwise idle. The book's own code is fast: the book update takes about
\qty{160}{\nano\second} at the median, the strategy about 100, the \omterm{def:ll:pre-trade-risk-and-kill-switches:gate}{risk gate} about 50 and the gateway about 30, all well under a
microsecond at the 99th percentile. The decode stage, about \qty{440}{\nano\second}, includes waiting for the earlier messages
of the same packet: five events arrive together and leave one after another. The ring's median of about \qty{0.8}{\micro\second} is the
same queueing, on the engine side, behind the orders of the previous events. The kernel's stages dwarf the rest: receiving a
UDP packet on the loopback interface takes about \qty{1.4}{\micro\second}, sending the order over TCP about \qty{2.9}{\micro\second},
and in the tail the ring waits milliseconds whenever the engine thread is preempted on its CPU (by the kernel, the hypervisor
or any other task),
which no user-space code can prevent on a machine that has not isolated it (chapter~13).

\begin{omfigure}
\begin{tikzpicture}
\begin{axis}[width=9cm,height=7.2cm,xbar,xmode=log,xmin=10,xmax=5e6,log origin=infty,bar width=3.2pt,
  ytick={0,1,2,3,4,5,6,7,8,9},yticklabels={receive,decode,ring,book,strategy,risk,gateway,transmit,tick to trade,in process},
  y dir=reverse,xlabel={nanoseconds},tick label style={font=\small},label style={font=\small},grid=major,
  grid style={gray!25},legend pos=outer north east,legend style={font=\scriptsize,cells={anchor=west}},table/col sep=comma]
\addplot[fill=omDef!55,draw=omDef] table[y expr=\coordindex,x=p50_ns]{figdata/low-latency/26-build-tick-to-trade-measured/measured_stages.csv};
\addlegendentry{median}
\addplot[fill=omMeth!45,draw=omMeth] table[y expr=\coordindex,x=p99_ns]{figdata/low-latency/26-build-tick-to-trade-measured/measured_stages.csv};
\addlegendentry{99th percentile}
\addplot[fill=omThm!40,draw=omThm] table[y expr=\coordindex,x=p999_ns]{figdata/low-latency/26-build-tick-to-trade-measured/measured_stages.csv};
\addlegendentry{99.9th percentile}
\end{axis}
\end{tikzpicture}

\omcaption{The assembled path, stage by stage and end to end: medians, 99th and 99.9th percentiles over ten runs (the stages on
the first order of each event, the totals on all 10\,500 orders), on a laptop (Intel Core Ultra 7 155H, WSL2) with no isolated
cores, the machine otherwise idle. Data: \texttt{bench\_t2t.py}.}\label{fig:ll:build-tick-to-trade-measured:stages}
\end{omfigure}

End to end, from the exchange's send of a packet to the venue's receipt of the order it caused, the median is about
\qty{9}{\micro\second}; in the process, from the packet's receipt to the order's message, about \qty{3.7}{\micro\second}
(\cref{fig:ll:build-tick-to-trade-measured:ccdf}). The distributions have the shape the book has predicted since chapter~1: a
tight body and a tail three orders of magnitude further out, set by the machine rather than the code. \Cref{tab:ll:build-tick-to-trade-measured:budget}
checks chapter~1's three-microsecond budget with \texttt{firm.latbudget}: the strategy and the \omterm{def:ll:pre-trade-risk-and-kill-switches:gate}{risk gate} meet their targets;
the book and the decode miss theirs (the book by a third, the decode because of the packet's queue), and the kernel's receive
and transmit miss by factors of 1.6 and 3.6, as they were bound to: the budget assumed a
kernel-bypass network card.

\begin{omfigure}
\begin{tikzpicture}
\begin{axis}[width=10cm,height=5.8cm,xmode=log,ymode=log,xmin=1000,xmax=5e6,ymin=0.001,ymax=1.1,
  xlabel={nanoseconds},ylabel={fraction of orders slower},tick label style={font=\small},label style={font=\small},
  grid=major,grid style={gray!25},legend pos=south west,legend style={font=\scriptsize,cells={anchor=west}},
  table/col sep=comma]
\addplot[omThm,thick,restrict expr to domain={\coordindex}{0:60}] table[x=ns,y=ccdf]{figdata/low-latency/26-build-tick-to-trade-measured/measured_t2t_ccdf.csv};
\addlegendentry{tick to trade (exchange to venue)}
\addplot[omDef,thick,restrict expr to domain={\coordindex}{61:121}] table[x=ns,y=ccdf]{figdata/low-latency/26-build-tick-to-trade-measured/measured_t2t_ccdf.csv};
\addlegendentry{in process (receipt to message)}
\end{axis}
\end{tikzpicture}

\omcaption{The end-to-end distributions of the assembled path, as the fraction of the 10\,500 orders slower than each latency
(log scales), on the same laptop. Data: \texttt{bench\_t2t.py}.}\label{fig:ll:build-tick-to-trade-measured:ccdf}
\end{omfigure}

\begin{table}[ht]
\centering\small
\begin{tabular}{lrrl}
\toprule
stage (median) & budget (ns) & measured (ns) & \\
\midrule
receive & 900 & 1\,400 & over \\
decode & 60 & 440 & over \\
book & 120 & 160 & over \\
strategy & 200 & 100 & within \\
risk & 50 & 47 & within \\
transmit & 800 & 2\,900 & over \\
end to end & 3\,000 & 7\,400 & over \\
\bottomrule
\end{tabular}
\caption{Chapter~1's three-microsecond budget checked against the measured stages with \texttt{firm.latbudget} (medians; the end
to end composed from the first orders' stages). Data: \texttt{bench\_t2t.py}.}\label{tab:ll:build-tick-to-trade-measured:budget}
\end{table}

\section{What a tuned server would change}

Three changes separate this laptop from a trading server, and each targets a part of the table. A kernel-bypass network card
removes the system calls from the data path: the \texttt{ef\_vi} interface of one vendor's adapters, for example, gives the application
``direct access'' to the adapter's datapath, with data path operations that ``do not require system calls''; chapter~1's budget
assumed \qty{0.9}{\micro\second} to receive and \qty{0.8}{\micro\second} to transmit on such a path, card and wire included. Isolated
cores (chapter~13) keep other work off the path's CPUs, which is what the tail of the ring stage is made of. And the path itself
can be shortened: the packet's queue behind the decode and the ring disappears if events are handed over as a batch or if the
feed thread runs the strategy for simple cases, and the book's \qty{160}{\nano\second} is the ladder of chapter~19 on a
small-tick instrument with a per-event engine call that the batch would amortise.

Replacing the measured receive and transmit by the budget's kernel-bypass figures and keeping the measured in-process stages
gives a projected median of about \qty{3.2}{\micro\second}: $0.9 + 0.8$ for the network, plus about \qty{1.5}{\micro\second} for
decode, ring, book, strategy, risk and gateway, of which the queueing in decode and ring is some two-thirds. That misses the three-microsecond
budget by the queue, which is the next thing to remove, and not by any of the components the book has built. The projection is
arithmetic on a stated assumption, not a measurement; the measurement on the target server is the only proof, and the harness
of this chapter is how it is taken.

\section{Tutorial: run the path, read the budget}\label{tut:ll:build-tick-to-trade-measured:tut}

\begin{tutorial}
\textbf{Goal.} Check the assembled path, run it live on the loopback interface, and read its stages against the budget.
\textbf{End state:} \cref{fig:ll:build-tick-to-trade-measured:stages,fig:ll:build-tick-to-trade-measured:ccdf}, the budget
table, and green tests in C++ and Rust.
\begin{tutsteps}
\item \textbf{The line.} \texttt{make\_t2t\_fixtures.py} encodes the small-tick day as the simulator's feed; the C++ test checks
that the \omterm{def:ll:the-feed-handler:handler}{feed handler} turns it back into the same 6\,000 events.
\item \textbf{Offline.} The C++ and Rust tests replay the events through the path and obtain the committed orders.
\item \textbf{Live.} \texttt{firm\_ticktotrade\_live\_test.cpp} runs the four threads once and checks the orders, the order of
the stamps and the absence of allocation; \texttt{python bench\_t2t.py} runs ten pinned runs and writes the stages, the
distributions and the budget report.
\end{tutsteps}
\textbf{What to change next.} Hand the events of a packet to the engine as one batch and measure the decode and ring stages
again; pin two threads on the two hyperthreads of one core and watch the stages that share it.
\end{tutorial}

\section{Build: tick-to-trade}\label{bld:ll:build-tick-to-trade-measured:ticktotrade}

\begin{build}
\bfield{Purpose} The book's components assembled into one trading path with its measurement harness: the proof that they compute
the right orders together, and the measurement of what each costs.
\bfield{Interface} C++20 \texttt{firm::t2t}: \texttt{Path(tick)} with \texttt{on\_event(event, now, stamp, send)}, \texttt{out},
\texttt{order\_hash}, \texttt{read\_events}, \texttt{Stamp}; \texttt{run\_live(Config)} returning \texttt{Result} with the
orders' stamps, the orders, the counts and the allocations after warm-up. Rust \texttt{firm\_ticktotrade}: \texttt{Path},
\texttt{events\_from\_line}, \texttt{order\_hash}. Python: \texttt{make\_t2t\_fixtures.py}; the chapter's \texttt{ll\_t2t.py}
computes the stages and checks the budget with \texttt{firm.latbudget}.
\bfield{Rules} The live path and the offline replay produce the same orders; every stamp is taken outside any loop, in path order;
no allocation after warm-up; the stages attributed on the first order of each event.
\bfield{Acceptance tests} \texttt{code/firm/ticktotrade/}: the recorded line decoded to the \omterm{def:ll:the-order-book-builder:builder}{book builder}'s events; the offline
orders equal to the committed hash in C++ and in Rust, with no allocation in a second pass; one live run with the same orders,
every order's stamps in order and received by the venue, and no allocation after warm-up.
\bfield{Stretch} The simulator's live server (Book 10) as the venue, with acknowledgements through the gateway's state machine;
a batched hand-over from the feed thread; hardware time stamps from a network card; the harness on a tuned server.
\end{build}

\begin{omsources}
\item Xilinx (AMD), \emph{ef\_vi User Guide}, SF-114063-CD, issue 10, 2019.
\end{omsources}

\section{Exercises}

\begin{exercise}[$\star$]\label{exo:ll:build-tick-to-trade-measured:1}
Using the medians of \cref{fig:ll:build-tick-to-trade-measured:stages}, what share of the in-process time is spent in the
book's own components (book, strategy, risk, gateway)?
\end{exercise}

\begin{exercise}[$\star$]\label{exo:ll:build-tick-to-trade-measured:2}
Why must the clock read at each \omterm{def:ll:build-tick-to-trade-measured:points}{instrumentation point} be the same counter on every core, and what would go wrong with a
per-thread clock?
\end{exercise}

\begin{exercise}[$\star$]\label{exo:ll:build-tick-to-trade-measured:3}
Five events arrive in one packet and each takes about \qty{100}{\nano\second} to decode and publish. What decode stage do the
five events see on average, and why does the ring stage show the same effect?
\end{exercise}

\begin{exercise}[$\star\star$]\label{exo:ll:build-tick-to-trade-measured:4}
An event causes two orders, each sent in about \qty{4}{\micro\second}. If the stages were attributed on every order, which stage
would absorb the wait of the second, and by how much would its median move?
\end{exercise}

\begin{exercise}[$\star\star$]\label{exo:ll:build-tick-to-trade-measured:5}
Recompute the projection of the last section if the queueing in decode and ring were removed entirely. Does the projected path
meet the three-microsecond budget?
\end{exercise}

\begin{exercise}[$\star\star$]\label{exo:ll:build-tick-to-trade-measured:6}
The receive stage includes the exchange thread's \texttt{sendto}. Why can this harness not separate the sender's system call from
the receiver's, and what would a hardware time stamp on the wire change?
\end{exercise}

\begin{exercise}[$\star\star\star$]\label{exo:ll:build-tick-to-trade-measured:7}
\emph{Coding.} Hand the events of a packet to the engine as one ring message and run the engine on the batch. Check that the
orders are unchanged and measure the decode and ring stages again.
\end{exercise}

\begin{exercise}[$\star\star\star$]\label{exo:ll:build-tick-to-trade-measured:8}
\emph{Find the flaw.} ``Our tick-to-trade is \qty{800}{\nano\second}: we time from the moment our code gets the packet to the
moment we call \texttt{send}, and we report the average.''
\end{exercise}

\section{Problem: The Budget, Proven}

\begin{problem}[{Weekend problem --- the path of the book, stage by stage}]\label{pb:ll:build-tick-to-trade-measured:1}
Use \cref{fig:ll:build-tick-to-trade-measured:stages,tab:ll:build-tick-to-trade-measured:budget} (\texttt{measured\_stages.csv},
\texttt{measured\_budget.csv}).

\medskip\noindent\textbf{Part I --- The measurement.}
\begin{enumerate}
\item List the eight stages and the \omterm{def:ll:build-tick-to-trade-measured:points}{instrumentation points} that bound each.
\item Give the median, 99th and 99.9th percentiles of the tick-to-trade and in-process latencies.
\item Which stage has the largest median, and which the largest 99.9th percentile?
\item Why is the 99.9th percentile of the ring stage so far from its median on this laptop?
\end{enumerate}

\medskip\noindent\textbf{Part II --- The check.}
\begin{enumerate}[resume]
\item Which stages meet chapter~1's budget at the median, and which miss it, by what factor?
\item Why is the budget's end-to-end figure composed from the stages of the first orders rather than taken from the
tick-to-trade distribution?
\item Why do the percentiles of the stages not add up to the percentile of the total?
\item What does the check establish that the functional tests do not?
\end{enumerate}

\medskip\noindent\textbf{Part III --- The projection.}
\begin{enumerate}[resume]
\item Replace the receive and transmit stages by the budget's kernel-bypass figures and compute the projected median.
\item Which part of the remaining gap is queueing, and which is computation?
\item What would isolated cores change in the figure, and what would they not change?
\item Why is the projection not a measurement, and what would make it one?
\end{enumerate}

\medskip\noindent\textbf{Part IV --- The verdict.}
\begin{enumerate}[resume]
\item State the \emph{named result}: the measured median, 99th and 99.9th percentiles per stage and end to end on this laptop,
and the projected median for a tuned server.
\item How much of the in-process median is the code this book wrote, and how much is waiting?
\item Which chapter's technique addresses each of the three largest contributions?
\item Why do the offline and live paths have to produce the same orders before any timing is trusted?
\item Why is the absence of allocation checked in every run?
\item What would you measure first on a production server?
\item Which result of the book would you distrust most if this harness had not been built?
\item In one sentence: what does it take to know a latency, rather than to believe it?
\end{enumerate}
\end{problem}

\section{Interview questions}

\begin{interviewq}[$\star$ \iqroles{developer}]\label{iq:ll:build-tick-to-trade-measured:1}
What is \omterm{def:ll:where-latency-comes-from:t2t}{tick-to-trade latency}, and where exactly do you start and stop the clock?
\end{interviewq}

\begin{interviewq}[$\star\star$ \iqroles{developer}]\label{iq:ll:build-tick-to-trade-measured:2}
How would you instrument a trading path to find out which stage is slow, without making it slower?
\end{interviewq}

\begin{interviewq}[$\star\star$ \iqroles{developer}]\label{iq:ll:build-tick-to-trade-measured:3}
Your median tick-to-trade is fine but the 99.9th percentile is a thousand times larger. Where do you look?
\end{interviewq}

\begin{interviewq}[$\star\star$ \iqroles{developer}]\label{iq:ll:build-tick-to-trade-measured:4}
Walk through a packet's path from the network card to your strategy and back out, with a latency for each step on a tuned
server.
\end{interviewq}

\begin{interviewq}[$\star\star$ \iqroles{developer, trader}]\label{iq:ll:build-tick-to-trade-measured:5}
The strategy team asks for a \omterm{def:ll:where-latency-comes-from:budget}{latency budget}. How do you write it, and how do you prove it is met?
\end{interviewq}

\begin{interviewq}[$\star\star\star$ \iqroles{developer}]\label{iq:ll:build-tick-to-trade-measured:6}
Design the measurement harness for a production trading system: what is stamped, where, with which clock, how the stamps are
joined, and how the results reach the people who need them.
\end{interviewq}
